\section{Source audit, proposed corrections, and further research}
\label{audit:sec}

\subsection{What the audit establishes}

The repository was inspected at the revision on the title page.
All six incoming ZIPs were unpacked with a path-safety check; their
inventories, claim/status files, and relevant sections were read.
The accompanying \texttt{incoming\_audit\_summary.json} records the
archive names, SHA-256 digests, already covered results, and open
questions.  This was a targeted mathematical audit of the claims
used here, not a line-by-line certification of every theorem in the
large source corpus.

The source claims used most directly are the exact even-tail
Bernoulli coordinates, the first-$q$ conjecture, the Gamma-quotient
ordinary integral, the fifth-moment baseline, and the recorded
Family--2 example.  The tail coefficient/operator conversion was
checked independently; the new proof includes the column reversal
and the omitted lower-index-zero convention.  The Gamma bridge
was reproved with its Laurent normalization before it was used in
the sixth moment.  The Family--2 equation was compared to its
original right-hand side, not merely to a numerically similar
simplification.

\textbf{No new false theorem in the inherited material was established
by this audit.}  Some incoming reports already correct signs,
Stieltjes indexing, or expressibility claims in external papers.
Those corrections are identified as inherited in the audit record;
they are not counted as new findings here.

\subsection{Concrete proposed source changes}

\begin{enumerate}
\item \textbf{Upgrade the Family--2 status.}
In the canonical \texttt{chapters/04-depth.tex}, the equation labeled
\texttt{double:eq:fam2-example} can be marked proved.
Retain its original identity for continuity, add
\eqref{bip:eq:family2-short} as its simplified form, and cite the
two exact component evaluations \eqref{bip:eq:12}--\eqref{bip:eq:21}.
The discovery by integer relation remains a useful historical note,
but it is no longer the justification for that equation.

\item \textbf{Refine the optimal-prefix conjecture.}
In the incoming canonical fifth-moment report, retain
\texttt{conj:firstq} for general unequal even orders.
Record $p=q-2$ as a theorem with determinant
\eqref{adj:determinant}; retain $p=q$ as the already proved source
case and general gaps $q-p\ge4$ as open.
This is an infinite-family resolution, not a complete proof of the
original conjecture.

\item \textbf{Attach the sixth calculation to the higher-moment agenda.}
The new result \eqref{sixth:Q6} completes a specific sixth-power
calculation in explicit convergent zeta-derivative coordinates.
Its remaining arithmetic question is the reduction of $\Omega_5$.
The generic all-power algorithm and the fifth calculation were
already present and should retain their earlier attribution.

\item \textbf{Cite the exact antecedents for two reusable methods.}
The functional-monomial direct sum should cite
\cite[Proposition~4.10]{Guo2014}; the bicolored parity implementation
should cite Nakamura and Panzer \cite{bip:Nakamura2012,bip:Panzer2017}.
The latter paper's footnote~16 identifies a transposed product in
Nakamura's printed equation (1.3).  Our proof does not import that
formula without rederivation.  This is an attribution and formula-use
clarification, not an assertion that the supplied manuscript copied
the misprint.

\item \textbf{Preserve the precise meanings of ``finite part'' and ``rank''.}
When importing the new results, keep the original-coordinate
definition \eqref{scope:finitepart}, the local correction in the
Hurwitz product section, and the formal coefficient-space qualifiers.
A raw Laurent constant of the spectral kernel generally differs
from the coordinate integral; $\NN$ is not multiplicative; and a
formal rank computation does not prove arithmetic independence.
These are explicit integration requirements for this article, not
newly alleged errors in every prior use of those terms.
\end{enumerate}

\subsection{Research questions with concrete mathematical targets}

The following agenda separates exact next calculations from genuine
structural questions.  An answer should supply a proof, a fully stated
functional identity, or an exact algebraic certificate.  A large
precision search or a numerical rank alone would not settle the
corresponding question.

\paragraph{1. The first unresolved even-tail gap.}
Prove or disprove the optimal first-$q$ conjecture for
$q=p+4$, uniformly in $p$.  The adjacent-order proof uses one free
formal coefficient and a sign-definite scalar recurrence.  The next
gap produces a larger initial-data space.  The concrete target is a
nonzero determinant formula for the associated finite block, perhaps
through a matrix recurrence whose relevant minors have a controlled
sign.  This would test whether the adjacent argument extends or
depends essentially on second order.

\paragraph{2. A determinant formula for arbitrary unequal orders.}
For $p=2a,q=2b$, seek a closed factorization or a finite positive
coefficient formula for the late Bernoulli matrix determinant.
Any proposed expression must retain the source's omitted binomial
terms and the centered cutoff.  A proof should also give the inverse
map, thereby producing ordinary harmonic-tail identities in addition
to a rank statement.

\paragraph{3. Arithmetic structure of the adjacent determinant sequence.}
The rational numbers $c_n$ have both a Bernoulli recurrence and the
central-binomial representation \eqref{adj:detgenerator}.
Determine an exact denominator formula, a useful integer
renormalization, or a finite convolution identity for this sequence.
The target is an arithmetic identity or recurrence with transparent
integrality, rather than an asymptotic estimate.

\paragraph{4. The next full Gamma coefficient sector.}
Find the generating function for all functional monomials containing
$u_2^2$ and arbitrary powers of $u_0,u_1$, including any terms
coupled to $u_2^3$ by factor-degree parity.
The ideal argument used for the complete $u_0,u_1$ sector suggests
working modulo a precisely stated higher-derivative monomial ideal.
Its preservation by the normal projection must be proved from the
reduction rule; stability under $D$ by itself is insufficient.

\paragraph{5. Higher sectors and a global normal-form generating function.}
The centered generator contains only even derivative indices before
normal reduction.  Can that symmetry be turned into a closed
multivariable generating function for $\NN(\A_m)$, beyond the
square and digamma--trigamma sectors?  An exact coefficient-extraction
formula indexed by partitions would be useful even without a short
elementary expression.

\paragraph{6. Arithmetic reduction of the mixed sixth coordinate.}
Determine whether $\Omega_5$ can be expressed in a prescribed
algebra generated by the height-one $\lambda_d$, ordinary zeta
derivatives, and lower mixed coordinates already appearing in the
fifth moment.  Appendix~\ref{tables:omega-convergent} supplies a
finite list of absolutely convergent MZV germs of depth at most three.
This makes the question precise enough for a symbolic functional
relation search.  Failure of a finite search would not prove
nonmembership or independence.

\paragraph{7. Complex-order functional relations before differentiation.}
The preceding question is about derivatives in zeta orders, not only
values at integral orders.  A parity identity known only for integer
indices cannot be differentiated in an index without an analytic
extension.  Seek functional equations for the required complex-order
germs, prove their continuation and local regularity, and only then
differentiate.  This is a specific way to avoid an unjustified
``differentiate the MZV reduction'' step.

\paragraph{8. Minimal spectral depth of higher moments.}
The normal-form and residue theorems give an upper bound on the
depth needed for the derivative-containing remainder of $Q_m$.
Determine when that bound is attained after all proven functional
relations are used.  First complete explicit seventh- and eighth-weight
coordinate tables, then look for uniform cancellations.
Distinguish the dimension of differential-polynomial normal forms,
the rank of their harmonic residue maps, and the dimension of their
actual period values; these are different problems.

\paragraph{9. Generalized Stieltjes products at the missing derivative edge.}
The finite generator in \eqref{hprod:stieltjes-family} treats every
pair of positive argument-derivative orders.  Extend its manifestly
holomorphic coefficient form to cases with one or both derivative
orders zero, and prove compatibility with the incoming
Stieltjes-collision formulas.  The extra spectral poles at order one
must be subtracted jointly; blindly setting an integer parameter to
one in the positive-order formula does not establish this extension.

\paragraph{10. Exact cyclotomic identities in the complementary parity sector.}
The bicolored theorem computes real parts in odd weight and imaginary
parts in even weight.  The Gaussian $S_6$ and revised $S_8$ questions
are not resolved by that mechanism.  Seek a genuinely additional
relation: a proven depth-three reduction, an integral symmetry, or a
specialization of a functional equation with a larger argument set.
The precise definitions and corrected candidates in the current
source should be retained; already rejected coefficient lists should
not be recycled as conjectures.

\paragraph{11. Antiderivatives and a specified function class.}
The spectral theorem provides normally convergent anchored primitives
for every derivative-containing differential polynomial.
Which of these primitives reduce to a specified finite class of
Gamma, Barnes, polylogarithm, and zeta-order-derivative functions?
A useful first target is the primitive of $\psi^2\psi'^2$.
Any nonexistence claim must name the allowed class and prove an
obstruction within it; an unspecified notion of ``closed form''
does not define a mathematical problem.

\paragraph{12. Exact proof certificates for repository integration.}
Turn the polynomial reductions, low-weight shuffle identities, and
finite Bernoulli matrix conjugacies into machine-readable exact
certificates.  The current scripts already separate rational
identities from numerical diagnostics.  A next stage could check
those certificates with an independent implementation or a proof
assistant, while leaving analytic continuation and contour estimates
as explicit lemmas with their assumptions stated.

The most immediate next proof target is the $q-p=4$ tail determinant.
The most immediate next arithmetic target is the convergent
depth-three expression for $\Omega_5$.  Both have precise finite
input data in this article, and both would produce exact identities
of direct use to the existing manuscript.
